\subsection{5. The Automation Scripts: What They Do and When to Use
Them}\label{the-automation-scripts-what-they-do-and-when-to-use-them}

\subsubsection{\texorpdfstring{\texttt{new\_experiment.sh} ---
Scaffolding New
Experiments}{new\_experiment.sh --- Scaffolding New Experiments}}\label{new_experiment.sh-scaffolding-new-experiments}

\textbf{When to use it:} Only when adding an experiment \textbf{beyond
the 17 already scaffolded} (i.e., experiment 018 onwards).

\textbf{What it does:} 1. Copies the entire
\texttt{experiments/\textbackslash\{\}\_template/} directory to a new
folder. 2. Replaces all placeholder text
(\texttt{[PAPER\textbackslash\{\}\_TITLE]},
\texttt{[CODE\textbackslash\{\}\_URL]}, etc.) in the new folder's files.
3. Sets today's date in the
\texttt{REPLICATION\textbackslash\{\}\_LOG.md}.

\textbf{How to use it:}

\begin{minted}{bash}
./scripts/new_experiment.sh NNN "folder_name" "Paper Title" "https://code-url"
\end{minted}

\textbf{Real example:}

\begin{minted}{bash}
./scripts/new_experiment.sh 018 "smogn_regression" "SMOGN: a Pre-processing Approach for Imbalanced Regression" "https://github.com/nicktaonline/smogn"
\end{minted}

\textbf{After running it, you will see:}

\begin{minted}{text}
[DIR] Creating experiment: 018_smogn_regression
[OK] Experiment scaffolded at: /Users/joaopms/Documents/Projeto_Mestrado/experiments/018_smogn_regression

Next steps:
  1. Edit README.md — fill in metadata fields
  2. Clone original code into src/original/
  3. Set up environment and update requirements.txt
  4. Run: python3 scripts/update_readme.py
  5. Commit: git add experiments/018_smogn_regression/ ...
\end{minted}

\begin{quote}
\small {[}WARN{]} \textbf{Do NOT run this script} for experiments 001--017.
They are already scaffolded.
\end{quote}

\subsubsection{\texorpdfstring{\texttt{update\_readme.py} --- Refreshing
the
Dashboard}{update\_readme.py --- Refreshing the Dashboard}}\label{update_readme.py-refreshing-the-dashboard}

\textbf{When to use it:} Every time you change the status badge in any
experiment's \texttt{README.md}.

\textbf{What it does:} 1. Scans every folder under
\texttt{experiments/NNN\textbackslash\{\}\_*/} for a
\texttt{README.md}. 2. Reads the
\texttt{\textbackslash\{\}\#\textbackslash\{\}\# Replication Status: [badge]} line from
each one. 3. Updates the progress summary table (Done / In Progress /
Setup / Pending / Blocked counts). 4. Updates today's date in the ``Last
updated'' field.

\textbf{How to use it:}

\begin{minted}{bash}
# Dry run — shows what WOULD change, but writes nothing
python3 scripts/update_readme.py --check

# Live run — actually updates README.md
python3 scripts/update_readme.py
\end{minted}

\textbf{When to run it:} - Whenever you change an experiment's status
badge - Before every \texttt{git commit} that involves
progress changes

\begin{quote}
\small \textbf{Important:} This script only updates the progress tables. It
does NOT change the Experiment Index table. If you add a brand new
experiment and want it in the index, you must manually add a row there.
\end{quote}

\subsubsection{\texorpdfstring{\texttt{sync\_from\_excel.py} --- Syncing
Paper
Metadata}{sync\_from\_excel.py --- Syncing Paper Metadata}}\label{sync_from_excel.py-syncing-paper-metadata}

\textbf{When to use it:} When you update the
\texttt{Fichamento\textbackslash\{\}\_Artigos.xlsx} file in Google Drive
and want the changes reflected in
\texttt{docs/paper\textbackslash\{\}\_registry.md}.

\textbf{What it does:} 1. Opens
\texttt{Fichamento\textbackslash\{\}\_Artigos.xlsx} from your Google
Drive (path is hardcoded in the script). 2. Reads all rows from the
first sheet. 3. Groups papers into ``with code'' and ``without code''
based on the \texttt{Disponibilidade de Códigos?}
column. 4. Writes a formatted markdown table to
\texttt{docs/paper\textbackslash\{\}\_registry.md}.

\textbf{How to use it:}

\begin{minted}{bash}
python3 scripts/sync_from_excel.py
\end{minted}

\textbf{Expected output:}

\begin{minted}{text}
[OK] Paper registry written to: .../docs/paper_registry.md
   17 papers with code, 26 without code
\end{minted}

\begin{quote}
\small {[}WARN{]} \textbf{Requirement:} \texttt{openpyxl} must
be installed (\texttt{pip3 install openpyxl}).
{[}WARN{]} \textbf{Requirement:} Google Drive must be mounted and the
file accessible at its hardcoded path.
\end{quote}
